\begin{figure}[htpb]
    \centering
    % Resizebox giúp sơ đồ tự động co giãn vừa vặn với chiều ngang trang giấy
    \resizebox{\textwidth}{!}{
    \begin{tikzpicture}[
        % Định nghĩa style cho các khối hình
        box/.style={rectangle, draw, thick, minimum width=2.8cm, minimum height=1.2cm, text width=2.5cm, align=center, fill=white, font=\small},
        dia/.style={diamond, draw, thick, minimum width=3cm, minimum height=2.5cm, text width=1.8cm, align=center, fill=white, font=\small},
        line/.style={draw, thick, -Stealth},
        labeltext/.style={font=\scriptsize, align=center}
    ]

        % 1. VẼ KHUNG SWIMLANE (LÀN ĐƯỜNG)
        % Khung ngoài cùng
        \draw[thick] (-9, 0.5) rectangle (9, 12);
        % Đường kẻ ngang tiêu đề
        \draw[thick] (-9, 11) -- (9, 11);
        % Đường kẻ dọc chia 2 làn
        \draw[thick] (0, 0.5) -- (0, 12);
        
        % Tiêu đề của 2 làn
        \node[font=\bfseries] at (-4.5, 11.5) {Hệ thống};
        \node[font=\bfseries] at (4.5, 11.5) {Cán bộ quản lý đô thị};

        % 2. ĐẶT CÁC KHỐI NODE VÀO VỊ TRÍ
        % Làn phải (Cán bộ)
        \node[dia] (loai_tac_vu) at (2, 9) {Loại tác vụ};
        \node[box] (nhap_moi) at (7, 9) {Nhập mới};
        \node[box] (cap_nhat) at (2, 5) {Cập nhật};
        \node[box] (nhap_thong_tin) at (7, 3.85) {Nhập thông tin mô tả};
        \node[box] (gui_yeu_cau) at (7, 2) {Gửi yêu cầu phê duyệt};

        % Làn trái (Hệ thống)
        \node[dia] (du_lieu_hop_le) at (-3, 9) {Dữ liệu hợp lệ};
        \node[box] (luu_tru) at (-7, 9) {Lưu trữ vào CSDL};
        \node[box] (yeu_cau_chinh_sua) at (-3, 3.5) {Yêu cầu chỉnh sửa nội dung};

        % 3. VẼ CÁC ĐƯỜNG MŨI TÊN LIÊN KẾT
        % Các mũi tên trong làn phải
        \draw[line] (loai_tac_vu.east) -- node[above, labeltext] {Nhập thông\\tin tài sản} (nhap_moi.west);
        \draw[line] (loai_tac_vu.south) -- node[right, labeltext] {Đo đạc thực\\địa/Cập nhật} (cap_nhat.north);
        \draw[line] (nhap_moi.south) -- (nhap_thong_tin.north);
        \draw[line] (cap_nhat.east) -- (nhap_thong_tin.west);
        \draw[line] (nhap_thong_tin.south) -- (gui_yeu_cau.north);

        % Mũi tên từ làn phải sang làn trái (Gửi yêu cầu -> Dữ liệu hợp lệ)
        % (Đi sang trái tới tọa độ X = -0.5, sau đó đi lên và rẽ trái vào khối)
        \draw[line] (gui_yeu_cau.west) -- (-0.5, 2) |- (du_lieu_hop_le.east);

        % Các mũi tên trong làn trái
        \draw[line] (du_lieu_hop_le.west) -- node[above, labeltext] {Đúng} (luu_tru.east);
        \draw[line] (du_lieu_hop_le.south) -- node[left, labeltext] {Từ chối} (yeu_cau_chinh_sua.north);

        % Mũi tên từ làn trái trả về làn phải (Yêu cầu chỉnh sửa -> Nhập thông tin)
        % (Đi sang phải tới tọa độ X = 1.2, sau đó đi lên và rẽ phải vào khối)
        % [yshift=-10pt] giúp mũi tên đi vào dưới đường của khối Cập nhật để không bị đè nét
        \draw[line] (yeu_cau_chinh_sua.east) -- (1.2, 3.5) |- ([yshift=-10pt]nhap_thong_tin.west);

    \end{tikzpicture}
    }
    \caption{Sơ đồ quy trình cập nhật thông tin tài sản}
    \label{fig:quy_trinh_cap_nhat}
\end{figure}